\documentclass[12pt]{article}
\usepackage[utf8]{inputenc}
\usepackage[T1]{fontenc}
\usepackage[margin=1in]{geometry}
\usepackage{setspace}
\usepackage{xcolor}
\usepackage[colorlinks=true, urlcolor=blue!60!black, linkcolor=blue!60!black]{hyperref}
\setlength{\parindent}{2em}
\setlength{\parskip}{1.25em}
\renewcommand{\baselinestretch}{1.0}
\begin{document}
\begin{center}
{\LARGE \textit{NAACP v. Alabama ex rel. Patterson}, 357 U.S. 449 (1958)\par}
\vspace{0.8em}
{\large Case Brief\par}
\end{center}
\vspace{1em}

\section{Memory Jogger}
Alabama tried to oust the NAACP and held it in contempt for refusing to hand over its rank-and-file membership lists; a unanimous Court held that compelled disclosure of members' identities, exposing them to reprisal without a sufficient justification supported by a substantial relation between disclosure and an important state interest, violates the freedom of association protected by the Fourteenth Amendment (\href{https://supreme.justia.com/cases/federal/us/357/449/}{Justia}).

\section{Detailed Case Facts}
Alabama's foreign-corporation qualification statute (Ala.\ Code, 1940, Tit.\ 10, \S\S\,192--198) required out-of-state corporations to file their corporate charter and designate a place of business and agent for service of process before transacting intrastate business. The National Association for the Advancement of Colored People (NAACP), a New York nonprofit membership corporation, had operated in Alabama through chartered local affiliates since 1918 and opened a regional office in 1951 without qualifying, believing itself exempt.

\section{Procedural History}
In 1956, Alabama Attorney General John Patterson filed an equity suit in Montgomery County Circuit Court to enjoin the NAACP from operating in Alabama and permanently oust it from the State. The complaint targeted the NAACP's civil rights advocacy, specifically citing its financial and legal support for Black students seeking admission to the University of Alabama and its support for the Montgomery bus boycott. On the day the suit was filed, the Circuit Court issued an \textit{ex parte} restraining order barring the NAACP both from conducting activities in Alabama and from taking any steps to qualify under state law.

When the NAACP moved to dissolve the restraining order, the State moved for production of extensive corporate records, including the names and addresses of all Alabama ``members'' and ``agents,'' claiming the records were needed to prove the NAACP was transacting intrastate business. The court ordered production. The NAACP admitted its Alabama activities, offered to qualify (submitting executed forms), and produced substantially all demanded records---including its charter, business records, names of directors and officers, total Alabama membership count, and dues totals---but refused on constitutional grounds to disclose its rank-and-file membership lists.

The Circuit Court adjudged the NAACP in civil contempt and imposed a \$10,000 fine, which automatically escalated to \$100,000 after five days of noncompliance. Under Alabama law, the contempt judgment also barred the NAACP from obtaining a merits hearing on the ouster suit until it purged the contempt. The Alabama Supreme Court twice dismissed petitions for certiorari (265 Ala.\ 349, 91 So.\ 2d 214), asserting that mandamus prior to contempt was the sole proper remedy. The U.S. Supreme Court granted certiorari.

\section{Judicial Votes}
The Supreme Court unanimously \textbf{reversed and remanded}, 9--0. The case was argued January 15--16, 1958, and decided June 30, 1958 (\href{https://supreme.justia.com/cases/federal/us/357/449/}{Justia}):

\begin{itemize}
    \item \textbf{Majority Opinion:} Justice Harlan delivered the opinion of a unanimous Court, joined by Chief Justice Warren and Justices Black, Frankfurter, Douglas, Burton, Clark, Brennan, and Whittaker.
    \item \textbf{Concurring Opinions:} None.
    \item \textbf{Dissenting Opinions:} None.
\end{itemize}

\section{Holding}
\textbf{Questions Presented:}

\begin{enumerate}
  \item \textbf{Jurisdiction and Standing:} Did the Alabama Supreme Court's procedural ruling constitute an adequate and independent state ground barring review, and did the NAACP have standing to assert the constitutional privacy rights of its nonparty rank-and-file members?
  \item \textbf{Freedom of Association:} Did Alabama's court order compelling disclosure of the NAACP's rank-and-file membership lists---and the resulting \$100,000 contempt fine---violate the First and Fourteenth Amendments?
\end{enumerate}

\textbf{Holding:} Reversed and remanded (9--0). First, the Alabama Supreme Court's novel procedural bar was not an adequate state ground, and the NAACP had associational standing to assert its members' constitutional rights. Second, freedom to engage in expressive association for the advancement of beliefs and ideas is an inseparable aspect of ``liberty'' under the Fourteenth Amendment's Due Process Clause (incorporating First Amendment protections). Compelled disclosure of rank-and-file membership lists violated the Constitution because disclosure exposed members to severe private reprisals and Alabama failed to demonstrate a compelling, closely related state interest. The contempt judgment and \$100,000 fine therefore fell.

\section{Analysis of Opinions}
\subsection{Majority Opinion (Justice Harlan)}

Writing for a unanimous Court, Justice Harlan structured the opinion in four steps:

\begin{enumerate}
  \item \textbf{No Adequate State Procedural Ground:} Reviewing Alabama precedent back to 1909 (including \textit{Ex parte Morris} (1949), a contempt review involving Ku Klux Klan membership lists), the Court found that Alabama had consistently reviewed contempt orders via certiorari and had even told the NAACP in an earlier stay denial that certiorari was the ``established'' method. ``Novelty in procedural requirements cannot be permitted to thwart review in this Court'' by litigants relying on prior state decisions (357 U.S.\ at 457--458).

  \item \textbf{Associational Standing:} Although litigants generally must assert their own rights, the NAACP could assert its rank-and-file members' anonymity rights because requiring members to sue in their own names ``would result in nullification of the right at the very moment of its assertion'' (357 U.S.\ at 459). Moreover, the Association and its members were ``in every practical sense, identical,'' and compelled disclosure directly threatened the NAACP's own membership and financial support.

  \item \textbf{Expressive Association, Privacy, and State Action:} Recognizing the close nexus between speech and assembly, Harlan held that freedom of association for advancing political, economic, religious, or cultural beliefs is protected by the Fourteenth Amendment and that state action curtailing it faces ``the closest scrutiny'' (357 U.S.\ at 460--461). Even indirect burdens violate the Constitution if they effectively deter association. Because uncontroverted evidence showed that past disclosure of NAACP members' identities exposed them to ``economic reprisal, loss of employment, threat of physical coercion, and other manifestations of public hostility,'' compulsory disclosure constituted a substantial restraint (357 U.S.\ at 462). Harlan rejected Alabama's defense that reprisals came from private citizens rather than the State: ``The crucial factor is the interplay of governmental and private action,'' since private retaliation takes hold only after state compulsion exposes the lists (357 U.S.\ at 463).

  \item \textbf{Failure of Compelling Interest and Tailoring:} Any subordinating state interest justifying such deterrence must be ``compelling'' (357 U.S.\ at 463). Alabama's sole stated purpose---determining whether the NAACP was conducting intrastate business under the foreign-corporation statute---bore no ``substantial bearing'' on rank-and-file names, because the NAACP had already admitted its Alabama presence since 1918, offered to qualify, and disclosed its charter, officers, business records, total membership, and dues. Finally, Harlan distinguished \textit{New York ex rel. Bryant v. Zimmerman}, 278 U.S.\ 63 (1928) (which upheld mandatory disclosure of Ku Klux Klan rosters), based on the Klan's history of unlawful intimidation and violence and its refusal to provide \textit{any} information about its local operations.
\end{enumerate}

\subsection{Subsequent Developments}

\begin{itemize}

\item \textbf{Civil Rights Progeny and Balancing:} The Court swiftly applied \textit{NAACP v. Alabama} to strike down municipal membership-disclosure ordinances (\textit{Bates v. City of Little Rock}, 361 U.S.\ 516 (1960)) and teacher association-disclosure laws (\textit{Shelton v. Tucker}, 364 U.S.\ 479 (1960)). By contrast, in campaign finance (\textit{Buckley v. Valeo}, 424 U.S.\ 1 (1976)) and referendum petitions (\textit{Doe v. Reed}, 561 U.S.\ 186 (2010)), the Court upheld general public-disclosure regimes under ``exacting scrutiny'' while preserving as-applied exemptions for groups showing a reasonable probability of threats or harassment.
  \item \textbf{\textit{Americans for Prosperity Foundation v. Bonta}, 594 U.S.\ 595 (2021)} (\href{https://supreme.justia.com/cases/federal/us/594/19-251/}{Justia}): Marking a major modern expansion of \textit{NAACP v. Alabama}, a 6--3 Court (Roberts, C.J.) facially invalidated California's requirement that charities file IRS Schedule~B major-donor lists confidentially with the State Attorney General. Clarifying ``exacting scrutiny,'' the Court held that compelled disclosure regimes must be \textit{narrowly tailored} to a sufficiently important government interest (even if not the least restrictive means), and that even non-public government collection creates an unconstitutional chilling risk when cast as a blanket dragnet for administrative convenience rather than targeted fraud investigation.

\item \textbf{\textit{First Choice Women's Resource Centers, Inc. v. Davenport}, 608 U.S.\ 174 (2026)} (\href{https://www.law.cornell.edu/supremecourt/text/24-781}{Cornell LII}; \href{https://www.supremecourt.gov/opinions/25pdf/608us1r30_3f14.pdf}{Slip Op.}): Decided unanimously on April 29, 2026 (Gorsuch, J.), the Court extended \textit{NAACP v. Alabama} and \textit{Bonta} to pre-enforcement Article~III standing. When New Jersey's Attorney General served a non-self-executing subpoena demanding a faith-based pregnancy center's donor identities, the Court held that the official demand \textit{itself} inflicted an immediate, ongoing First Amendment associational injury by chilling donors---even before a state court ordered enforcement and even if the State promised confidentiality.

\end{itemize}

\section{Examples: Future Applications}
\begin{enumerate}

\item \textbf{Same-Side (Disclosure Invalid): A Legislative Demand for an Advocacy Group's Members.} A state legislative committee investigating whether a small LGBTQ advocacy group in a rural county has lobbied without registering subpoenas the group's complete membership list. The group has already registered, produced its officers' names and its budget, and offered sworn testimony about its lobbying. It submits affidavits showing that members whose affiliation became known locally lost jobs and received threats. As in \textit{NAACP v.\ Alabama}, compelled disclosure would expose members to private reprisal made possible by state compulsion, and the rank-and-file names have no substantial bearing on the committee's stated purpose, which the group's other disclosures already serve. The subpoena cannot be enforced (\href{https://supreme.justia.com/cases/federal/us/357/449/}{Justia}).

\item \textbf{Same-Side (Disclosure Invalid): Teachers Required to List Every Association.} A State requires every public-school teacher, as a condition of employment, to file an affidavit listing every organization to which the teacher has belonged or contributed in the past five years. The State says it needs the information to evaluate teachers' fitness. Applying \textit{NAACP v.\ Alabama}, the Court struck down a nearly identical Arkansas law in \textit{Shelton v.\ Tucker}, 364 U.S.\ 479 (1960). The State may inquire into teachers' competence, but a demand for every association, with no limit to those relevant to fitness, reaches far beyond that interest and deters teachers from joining unpopular groups. Because disclosure is not closely tailored to the State's legitimate need, the requirement violates the freedom of association.

\item \textbf{Opposite-Side (Disclosure Valid): Contributions to a Federal Candidate.} A federal statute requires every candidate's campaign committee to report the name and address of anyone who contributes more than \$200, and the reports are made public. A donors' group challenges the law on its face, invoking \textit{NAACP v.\ Alabama}, but offers no evidence that its supporters have faced threats or harassment. Under \textit{Buckley v.\ Valeo}, 424 U.S.\ 1 (1976), disclosure of contributions to candidates survives exacting scrutiny because it informs voters, deters corruption, and supports enforcement of contribution limits. \textit{NAACP v.\ Alabama} still matters at the margin: a minor party or group that shows a reasonable probability that disclosure will bring threats or reprisals can seek an as-applied exemption.

\item \textbf{Opposite-Side (Disclosure Valid): Signers of a Referendum Petition.} Voters sign a petition to put a contested ballot measure before the electorate, and state law makes the signed petitions available to the public on request. Sponsors of the petition sue to block release of the signers' names, arguing that disclosure will chill participation. Under \textit{Doe v.\ Reed}, 561 U.S.\ 186 (2010), a facial challenge fails: signing a referendum petition is a public act in the lawmaking process, and disclosure serves the State's interest in detecting fraud and preserving the integrity of the election. Unlike the NAACP's members, the signers have not shown that disclosure will expose them to reprisal; if a particular group can make that showing, an as-applied challenge remains available.

\item \textbf{Fence-Sitter (Confidential Donor Lists and Foreign Influence): The Civic Transparency Act.} To combat foreign disinformation in state policy debates, a State enacts the ``Civic Transparency Act,'' requiring any nonprofit that spends over \$5,000 annually on public-issue ads to file with the Secretary of State a confidential list of all donors contributing over \$250, subject to a \$50,000 daily contempt fine for noncompliance. The Attorney General serves a subpoena on a controversial environmental-advocacy group demanding its full donor roster to audit compliance. The group sues immediately in federal court. For the group: Under \textit{NAACP v. Alabama} and \textit{First Choice Women's Resource Centers v. Davenport} (2026), the group can allege a present injury to its own First Amendment rights because an outstanding subpoena for private donor information can deter donors and chill the group's advocacy even before enforcement. On the merits, under \textit{NAACP v. Alabama} and \textit{Americans for Prosperity Foundation v. Bonta} (2021), the law fails exacting scrutiny: even confidential government collection chills dissident association, and a blanket \$250 donor dragnet is not narrowly tailored to detecting foreign interference when the State could use targeted, cause-based subpoenas or require foreign-funding certifications without sweeping up domestic rank-and-file donors. For the State: Invoke \textit{Buckley v. Valeo} and \textit{Doe v. Reed} to argue that combating foreign influence in public policy debates is a compelling governmental interest and that statutory confidentiality shields donors from the private physical and economic reprisals documented in \textit{NAACP v. Alabama}. Distinguish \textit{Bonta} by arguing the statute targets only organizations actively spending on public-issue campaigns rather than every registered charity, making the \$250 threshold substantially related to enforcing foreign-funding bans. Because the law is aimed at a weighty interest but sweeps in every small domestic donor, its fate turns on how strictly a court applies \textit{Bonta}'s tailoring requirement.

\end{enumerate}

\section{Critique}
\begin{itemize}
  \item \textit{Unenumerated Right Grounded in First Amendment Architecture:} The word ``association'' does not appear in the First Amendment. Harlan derived expressive association functionally---as an indispensable amplifier of speech and assembly---and anchored it in Fourteenth Amendment ``liberty,'' laying analytical groundwork later invoked in privacy cases like \textit{Griswold v. Connecticut}, 381 U.S.\ 479 (1965).
  \item \textit{From Civil Rights Shield to Deregulatory Sword?} Critics (including Justice Sotomayor's dissent in \textit{Bonta}) argue that \textit{NAACP v. Alabama} turned on an extraordinary, uncontroverted record of Jim Crow violence and pretextual targeting, whereas applying facial exacting scrutiny to routine charitable or political donor disclosures risks insulating ``dark money'' from legitimate transparency and anti-fraud oversight.
\end{itemize}

\section{Key Quotations}
\begin{quote}
``To require that it be claimed by the members themselves would result in nullification of the right at the very moment of its assertion.''
\end{quote}
\noindent --- Justice Harlan, 357 U.S.\ at 459.

\begin{quote}
``It is beyond debate that freedom to engage in association for the advancement of beliefs and ideas is an inseparable aspect of the `liberty' assured by the Due Process Clause of the Fourteenth Amendment, which embraces freedom of speech. \dots\ [S]tate action which may have the effect of curtailing the freedom to associate is subject to the closest scrutiny.''
\end{quote}
\noindent --- Justice Harlan, 357 U.S.\ at 460--461.

\begin{quote}
``Inviolability of privacy in group association may in many circumstances be indispensable to preservation of freedom of association, particularly where a group espouses dissident beliefs.''
\end{quote}
\noindent --- Justice Harlan, 357 U.S.\ at 462.

\begin{quote}
``The crucial factor is the interplay of governmental and private action, for it is only after the initial exertion of state power represented by the production order that private action takes hold.''
\end{quote}
\noindent --- Justice Harlan, 357 U.S.\ at 463.


\end{document}
